% Part: first-order-logic
% Chapter: syntax-and-semantics
% Section: first-order-languages

\documentclass[../../../include/open-logic-section]{subfiles}

\begin{document}

\olfileid{fol}{syn}{fol}

\olsection{మొదటిస్థాయి విధేయ భాషలు}


మొదటిస్థాయి విధేయ తర్కపు వ్యక్తీకరణలను \emph{\tetoken{చరాలు}{!!{variable}s}},
\emph{\tetoken{స్థిరసంకేతాలు}{!!{constant}s}}, \emph{\tetoken{విధేయ సంకేతాలు}{!!{predicate}s}}, కొన్నిసార్లు
\emph{\tetoken{ప్రమేయ సంకేతాలు}{!!{function}s}} ఉన్న ప్రాథమిక పదజాలం నుంచి నిర్మిస్తాం. వాటితో
పాటు తార్కిక సంయోజకాలు, పరిమాణీకరణికాలు, కుండలీకరణ చిహ్నాలు, కామాలు
వంటి విరామసంకేతాలను ఉపయోగించి \emph{పదాలు}, \emph{\tetoken{సూత్రాలు}{!!{formula}s}}
ఏర్పరుస్తాం.

\begin{explain}
అనౌపచారికంగా, \tetoken{విధేయ సంకేతాలు}{!!{predicate}s} ధర్మాలు, సంబంధాలకు పేర్లు;
\tetoken{స్థిరసంకేతాలు}{!!{constant}s} వ్యక్తిగత వస్తువులకు పేర్లు; \tetoken{ప్రమేయ సంకేతాలు}{!!{function}s}
అన్వయాలకు పేర్లు. \tetoken{తాదాత్మ్య విధేయం}{!!{identity}}~$\eq$ మినహా ఇవన్నీ
\emph{తార్కికేతర సంకేతాలు}; ఇవన్నీ కలిసి ఒక భాషను ఏర్పరుస్తాయి.
ఏ మొదటిస్థాయి విధేయ భాష~$\Lang L$ అయినా దాని తార్కికేతర సంకేతాల ద్వారా
నిర్ణయితమవుతుంది. అత్యంత సాధారణ సందర్భంలో, $\Lang L$లో ప్రతి రకానికి
అనంతంగా అనేక సంకేతాలు ఉంటాయి.
\end{explain}

సాధారణ సందర్భంలో మొదటిస్థాయి విధేయ తర్కంలో కింది సంకేతాలను వాడతాం:

\begin{enumerate}
\item తార్కిక సంకేతాలు
\begin{enumerate}
\item తార్కిక సంయోజకాలు:
  \startycommalist
  \iftag{prvNot}{\ycomma $\lnot$ (నిషేధం)}{}%
  \iftag{prvAnd}{\ycomma $\land$ (సంధానం)}{}%
  \iftag{prvOr}{\ycomma $\lor$ (వికల్పం)}{}%
  \iftag{prvIf}{\ycomma $\lif$ (\tetoken{సోపాధికం}{!!{conditional}})}{}%
  \iftag{prvIff}{\ycomma $\liff$ (\tetoken{ద్విసోపాధికం}{!!{biconditional}})}{}%
  \iftag{prvAll}{\ycomma $\lforall$ (సార్వత్రిక పరిమాణీకరణికం)}{}%
  \iftag{prvEx}{\ycomma $\lexists$ (అస్తిత్వ పరిమాణీకరణికం)}{}.
\tagitem{prvFalse}{\tetoken{అసత్యానికి}{!!{falsity}} ప్రతిజ్ఞావాక్య స్థిరసంకేతం~$\lfalse$.}{}
\tagitem{prvTrue}{\tetoken{సత్యానికి}{!!{truth}} ప్రతిజ్ఞావాక్య స్థిరసంకేతం~$\ltrue$.}{}
\item ద్విస్థానిక \tetoken{తాదాత్మ్య విధేయం}{!!{identity}}~$\eq$.
\item \tetoken{చరాలు}{!!{variable}s} యొక్క \tetoken{అనంతంగా లెక్కించదగిన}{!!{denumerable}s} సమితి: $\Obj v_0$, $\Obj v_1$, $\Obj
  v_2$, \dots
\end{enumerate}
\item మొదటిస్థాయి విధేయ తర్కపు \emph{ప్రామాణిక భాష}ను ఏర్పరచే
  తార్కికేతర సంకేతాలు
\begin{enumerate}
\item ప్రతి $n>0$కు $n$-స్థాన \tetoken{విధేయ సంకేతాలు}{!!{predicate}s} యొక్క \tetoken{అనంతంగా లెక్కించదగిన}{!!{denumerable}s}
  సమితి: $\Obj A^n_0$, $\Obj A^n_1$, $\Obj A^n_2$, \dots
\item \tetoken{స్థిరసంకేతాలు}{!!{constant}s} యొక్క \tetoken{అనంతంగా లెక్కించదగిన}{!!{denumerable}s} సమితి: $\Obj c_0$, $\Obj c_1$, $\Obj
  c_2$, \dots.
\item ప్రతి $n>0$కు $n$-స్థాన \tetoken{ప్రమేయ సంకేతాలు}{!!{function}s} యొక్క \tetoken{అనంతంగా లెక్కించదగిన}{!!{denumerable}s}
  సమితి: $\Obj f^n_0$, $\Obj f^n_1$, $\Obj f^n_2$, \dots
\end{enumerate}
\item విరామసంకేతాలు: (, ), కామా.
\end{enumerate}

మన నిర్వచనాలు, ఫలితాల్లో అధిక భాగాన్ని మొదటిస్థాయి విధేయ తర్కపు పూర్తి
ప్రామాణిక భాషకోసం రూపొందిస్తాం. అయితే అనువర్తనాన్ని బట్టి భాషను కొన్ని
\tetoken{విధేయ సంకేతాలు}{!!{predicate}s}, \tetoken{స్థిరసంకేతాలు}{!!{constant}s}, \tetoken{ప్రమేయ సంకేతాలకు}{!!{function}s} మాత్రమే పరిమితం చేయవచ్చు.

\begin{ex}
అంకగణిత భాష~$\Lang L_A$లో ఒక్క ద్విస్థానిక \tetoken{విధేయ సంకేతం}{!!{predicate}}~$<$, ఒక్క
\tetoken{స్థిరసంకేతం}{!!{constant}}~$\Obj 0$, ఒక్క ఏకస్థానిక \tetoken{ప్రమేయ సంకేతం}{!!{function}}~$\prime$, రెండు
ద్విస్థానిక \tetoken{ప్రమేయ సంకేతాలు}{!!{function}s}~$+$,~$\times$ ఉంటాయి.
\end{ex}

\begin{ex}
సమితి సిద్ధాంత భాష~$\Lang L_Z$లో ఒక్క ద్విస్థానిక \tetoken{విధేయ సంకేతం}{!!{predicate}}~$\in$
మాత్రమే ఉంటుంది.
\end{ex}

\begin{ex}
క్రమాల భాష~$\Lang L_\le$లో ద్విస్థానిక \tetoken{విధేయ సంకేతం}{!!{predicate}}~$\le$ మాత్రమే
ఉంటుంది.
\end{ex}

మళ్ళీ చెప్పాలంటే ఇవి సంప్రదాయాలు మాత్రమే: అధికారికంగా ఇవి కేవలం
మారుపేర్లు. ఉదాహరణకు $<$, $\in$, $\le$ అనేవి $\Obj A^2_0$కు,
$\Obj 0$ అనేది $\Obj c_0$కు, $\prime$ అనేది $\Obj f^1_0$కు, $+$ అనేది
$\Obj f^2_0$కు, $\times$ అనేది $\Obj f^2_1$కు మారుపేర్లు.

\iftag{defNot,defOr,defAnd,defIf,defIff,defTrue,defFalse,defEx,defAll}{%
పైన ప్రవేశపెట్టిన ప్రాథమిక సంయోజకాలు,
\iftag{notprvEx,notprvAll}{పరిమాణీకరణికానికి}{పరిమాణీకరణికాలకు} అదనంగా,
కింది \emph{నిర్వచిత} సంకేతాలను కూడా వాడతాం:
\startycommalist
  \iftag{defNot}{\ycomma $\lnot$ (నిషేధం)}{}%
  \iftag{defAnd}{\ycomma $\land$ (సంధానం)}{}%
  \iftag{defOr}{\ycomma $\lor$ (వికల్పం)}{}%
  \iftag{defIf}{\ycomma $\lif$ (\tetoken{సోపాధికం}{!!{conditional}})}{}%
  \iftag{defIff}{\ycomma $\liff$ (\tetoken{ద్విసోపాధికం}{!!{biconditional}})}{}%
  \iftag{defAll}{\ycomma $\lforall$ (సార్వత్రిక పరిమాణీకరణికం)}{}%
  \iftag{defEx}{\ycomma $\lexists$ (అస్తిత్వ పరిమాణీకరణికం)}{}%
  \iftag{defFalse}{\ycomma \tetoken{అసత్యం}{!!{falsity}}~$\lfalse$}{}%
  \iftag{defTrue}{\ycomma \tetoken{సత్యం}{!!{truth}}~$\ltrue$}}{}.

\begin{tagblock}{defNot,defOr,defAnd,defIf,defIff,defTrue,defFalse,defEx,defAll}
\begin{explain}
నిర్వచిత సంకేతం అధికారికంగా భాషలో భాగం కాదు; దాన్ని అనౌపచారిక
సంక్షిప్తీకరణగా ప్రవేశపెడతాం. ప్రాథమిక సంకేతాలనే వాడితే చాలా పొడవయ్యే
సూత్రాలను సంక్షిప్తంగా రాయడానికి అది వీలు కల్పిస్తుంది. ఇది స్పష్టమైన
ప్రయోజనం. ఇంకా పెద్ద ప్రయోజనం ఏమిటంటే, నిరూపణలు చిన్నవవుతాయి.
ఒక సంకేతం ప్రాథమికమైతే నిరూపణల్లో దాని సందర్భాన్ని విడిగా
పరిశీలించాలి. అందువల్ల ప్రాథమిక సంకేతాలు ఎక్కువైతే మన నిరూపణలు కూడా
పొడవవుతాయి.
\end{explain}
\end{tagblock}

% Alternate symbols

\begin{intro}
మేము పైన వాడిన వాటికి భిన్నమైన పదజాలం, సంకేతాలు మీకు పరిచయమై
ఉండవచ్చు. తర్కశాస్త్ర గ్రంథాలు (బోధకులు కూడా) ``నిషేధం'' కోసం సాధారణంగా
$\sim$, $\neg$, లేదా~!ను, ``సంధానం'' కోసం $\wedge$, $\cdot$, లేదా $\&$ను
వాడతాయి. ``సోపాధికం'' లేదా ``ఆపాదన''కు సాధారణ సంకేతాలు
$\rightarrow$, $\Rightarrow$, $\supset$.
\iftag{prvIff,defIff}{``ద్విసోపాధికం,'' ``ద్వి-ఆపాదన,'' లేదా
  ``(భౌతిక) తుల్యత''కు సంకేతాలు $\leftrightarrow$,
  $\Leftrightarrow$, $\equiv$.}{}
\iftag{prvFalse,defFalse}{$\lfalse$ సంకేతాన్ని ``అసత్యం,''
  ``ఫాల్సమ్,'' ``అసంభవం,'' లేదా ``అధస్తనం'' అని పలు విధాలుగా అంటారు.}{}
\iftag{prvTrue,defTrue}{$\ltrue$ సంకేతాన్ని ``సత్యం,'' ``వేరమ్,'' లేదా
  ``ఉపరిస్థానం'' అని పలు విధాలుగా అంటారు.}{}

\tetoken{స్థిరసంకేతాలకు}{!!{constant}s} (వీటిని కొన్నిసార్లు పేర్లు అంటారు) లాటిన్ వర్ణమాల
మొదట్లోని చిన్న అక్షరాలు (ఉదా., $a$, $b$, $c$), \tetoken{చరాలకు}{!!{variable}s}
చివర్లోని చిన్న అక్షరాలు (ఉదా., $x$, $y$, $z$) వాడడం సంప్రదాయం.
పరిమాణీకరణికాలు \tetoken{చరాలతో}{!!{variable}s}, ఉదా., $x$తో కలుస్తాయి. సార్వత్రిక
పరిమాణీకరణికానికి $\forall x$, $(\forall x)$, $(x)$, $\Pi x$,
$\bigwedge_x$; అస్తిత్వ పరిమాణీకరణికానికి $\exists x$, $(\exists
x)$, $(Ex)$, $\Sigma x$, $\bigvee_x$ అనే సంకేతన భేదాలు ఉన్నాయి.
\end{intro}

\begin{explain}
ప్రతిజ్ఞావాక్య సంయోజకాలన్నింటినీ, రెండు పరిమాణీకరణికాలనూ భాషలో
ప్రాథమిక సంకేతాలుగా పరిగణించవచ్చు. లేదా ప్రాథమిక సంకేతాల చిన్న
సముదాయాన్ని ఎంచుకొని, మిగతా \tetoken{తార్కిక సంయోజకాలను}{!!{operator}s} నిర్వచితమైనవిగా
పరిగణించవచ్చు. బూలియన్ సంయోజకాల ``సత్య-ప్రమేయాత్మకంగా సంపూర్ణమైన''
సమితుల్లో $\{ \lnot, \lor \}$, $\{ \lnot, \land \}$, $\{ \lnot,
\lif\}$ ఉన్నాయి---వ్యక్తీకరణపరంగా సంపూర్ణమైన మొదటిస్థాయి విధేయ భాష
కోసం వీటిలో దేనినైనా ఏదో ఒక పరిమాణీకరణికంతో కలపవచ్చు.

మరో రెండు \tetoken{తార్కిక సంయోజకాలు}{!!{operator}s} మీకు పరిచయమై ఉండవచ్చు: హెన్రీ షెఫర్ పేరు
మీద వచ్చిన షెఫర్ రేఖ~$|$, క్వైన్ కటారి అని కూడా పిలిచే పియర్స్
బాణం~$\downarrow$. వీటికి వరుసగా ``న్యాండ్'', ``నార్'' అనే సాధారణ
పఠనాలను ఇచ్చినప్పుడు, ఒక్కొక్క సంయోజకమే సత్య-ప్రమేయాత్మకంగా
సంపూర్ణమైనది.
\end{explain}

\end{document}
